\documentclass[a4paper,12pt]{article}
\usepackage[utf8]{inputenc}
\usepackage[russian]{babel}
\usepackage{amsmath, amssymb}

\begin{document}

\section*{Задача 1}

Рассмотрим мультипликативную группу поля $GF(5^2)$. Она является циклической группой порядка $24$ и порождается некоторым примитивным элементом $x$.

Множество элементов поля можно записать в виде:
\[
\{0\} \cup \{x^0, x^1, \dots, x^{23}\}.
\]

Порядок нулевого элемента в мультипликативной группе не определён. Для любого ненулевого элемента $x^k$ его порядок вычисляется по формуле:
\[
\operatorname{ord}(x^k) = \frac{24}{\gcd(k, 24)}.
\]

Результаты вычислений представлены в таблице:

\[
\begin{array}{c|c}
k & \operatorname{ord}(x^k) \\
\hline
0 & 1 \\
1 & 24 \\
2 & 12 \\
3 & 8 \\
4 & 6 \\
5 & 24 \\
6 & 4 \\
7 & 24 \\
8 & 3 \\
9 & 8 \\
10 & 12 \\
11 & 24 \\
12 & 2 \\
13 & 24 \\
14 & 12 \\
15 & 8 \\
16 & 3 \\
17 & 24 \\
18 & 4 \\
19 & 24 \\
20 & 6 \\
21 & 8 \\
22 & 12 \\
23 & 24 \\
\end{array}
\]

Для нахождения порядков элементов использовалась следующая программа:

\begin{verbatim}
import math

for k in range(24):
    order = 24 // math.gcd(k, 24)
    print(k, order)
\end{verbatim}

\section*{Задача 2}

Рассматривается циклический код длины $15$ и размерности $4$ с проверочным полиномом
\[
h(x) = 1 + x + x^4
\]
над полем $\mathbb{F}_2$.

Код реализуется с помощью регистра сдвига с линейной обратной связью длины $4$. В качестве начального состояния выбрано:
\[
0011.
\]

Последовательность состояний регистра:
\[
0011 \to 1000 \to 0100 \to 0010 \to 0001 \to 1001 \to 1101 \to
\]
\[
1111 \to 1110 \to 0111 \to 1010 \to 0101 \to 1011 \to 1100 \to
\]
\[
0110 \to 0011.
\]

Регистр проходит через все $15$ ненулевых состояний и возвращается к исходному. Следовательно, период равен $15$, что соответствует максимальной длине ($2^4 - 1$).

Минимальное расстояние данного кода:
\[
d_{\min} = 8.
\]
Этот код является симплекс-кодом с параметрами $(15,4,8)$, который двойственен к коду Хэмминга $(15,11,3)$. Все ненулевые кодовые слова имеют одинаковый вес, равный $8$.

Для моделирования работы регистра и проверки периода использовался следующий код:

\begin{verbatim}
def simulate_lfsr(initial=[0,0,1,1]):
    state = initial[:]
    seen = set()

    for _ in range(20):
        tup = tuple(state)
        if tup in seen:
            break
        seen.add(tup)

        fb = state[0] ^ state[3]
        state = state[1:] + [fb]

    return len(seen)

print(simulate_lfsr())
\end{verbatim}

Для вычисления минимального расстояния был выполнен перебор всех кодовых слов:

\begin{verbatim}
import numpy as np

G = np.array([
[1,0,0,0,1,0,0,1,0,1,1,1,0,1,0],
[0,1,0,0,0,1,0,0,1,0,1,1,1,0,1],
[0,0,1,0,1,0,1,0,0,1,0,1,1,1,0],
[0,0,0,1,0,1,0,1,0,0,1,0,1,1,1]
], dtype=int) % 2

weights = []
for i in range(16):
    u = np.array([int(b) for b in format(i, '04b')])
    c = (u @ G) % 2
    weights.append(np.sum(c))

d_min = min(w for w in weights if w > 0)
print(d_min)
\end{verbatim}

\end{document}